\section{DeltaParq System Design}

\subsection{System Overview}

\subsubsection{Infrastructure}
For the design of the DeltaParq system, the structure and resources 
of the current system has been taken into account. The OPENIntel system 
utilizes a Hadoop file system (HDFS) running on top of an S3 storage cluster\cite{7460220}.
Furthermore, Apache Spark, with its Python derivative psypark, is used for 
querying and analyzing the dataset. While the DeltaParq system has been 
designed for the OPENIntel project, it has also been designed with 
generalizability in mind. The system uses a configuration file to set 
which columns are part of the primary key(PK) and which columns are volatile, 
this way other datasets could utilize the DeltaParq as well. All join operations 
within the system utilize this PK as the join condition. Furthermore, 
the integrations of the DeltaParq are extendible. Integrations with different 
file systems are easy to integrate and can after integration easily be switched 
in the configuration.


\subsubsection{Self-Describing Storage}

Like the current OPENIntel system, DeltaParq stores data using a Hive-style path on an S3 
bucket. DeltaParq constructs the path using the TLD, date, and S3 bucket that has been set 
in the configuration.
$$"s3a://{bucket}/tld={TLD}/year={year:04}/month={month:02}/day={day:02}/"$$
\todo[inline]{Fix this formatting}
Using this hierarchy, DeltaParq treats each month as a self-contained unit, this way each delta 
can only depend on a day in the same month. As the first day of the month is always the start, a 
part of the configuration called the $month\_config$ is stored in the month's directory. This part of  
the config describes the primary key, volatile columns, and strategy. This makes each stored month 
self-describing and ensures each dataframe can be correctly reconstructed without relying on the 
current system's configuration.

\subsubsection{Storage Strategies}

Since a delta will always be based on a single base dataframe, a storage strategy which determines 
which delta is based on which base, a strategy can be encoded as a function $strategy: \mathbb{Z} \rightarrow 
\mathbb{Z}$. Furthermore, a day is stored as a full snapshot if $strategy(d) = d$.

\subsection{Putting Data into the DeltaParq}

\subsubsection{Schema Decomposition}

As mentioned in the problem statement, the OPENIntel dataset 
contains both volatile and non-volatile data. As the volatile data is almost guaranteed 
to change between each measurement, these columns need to be removed from the non-volatile 
data using a schema decomposition. Using pyspark, the decomposition can be easily done by 
selecting only the necessary columns of a dataframe. For the volatile data, both the columns 
marked as volatile in the configuration and the columns of the primary key are selected 
such that this data can be joined with the non-volatile data for which all original columns, 
without the columns marked as volatile, but including the primary key, are selected. 
The resulting volatile dataframe can immediately be written to storage, while non-volatile 
dataframe needs to be processed using delta generation to be stored in the DeltaParq.

\subsubsection{Delta Generation}

The concept of delta generation is based on filtering out all the rows that are different 
in the updated version compared to the base version. The OPENIntel dataset, and DNS records 
in general, have an inherent flaw making this difficult: DNS records do not have a unique identifier.
The identifiers for a DNS record is its domain name and record type, but this combination does not have to be
unique. For example, the domain name example.com could return multiple TXT records, if between measurements
one TXT record changes while another is removed, there is no way of keeping track which of the TXT records 
changed and which one got removed. Not having a unique identifier leaves two options. Either every change is tracked 
as an addition or subtraction, this would result in two subtractions and one addition being stored in the above example, 
or the whole set of rows identified by the domain name and type is deemed as updated and thus stored in the delta.
For the DeltaParq, the latter option of storing the full identified block was chosen. \todo[inline]{explain why!!!}
A simple subtraction based algorithm, such as Algorithm~\ref{alg:subtract-based-delta} uses the pyspark subtract function 
to gather all rows in the updated dataframe not present in the removed dataframe. Rows that do not exist in the update 
but do exist in the base dataframe will not show up in this subtraction, thus a second subtraction needs to be done 
to gather all rows in the base dataframe not present in the update. Combining these keys together 
then results in all keys necessary for the delta.

\begin{algorithm}
    \caption{Subtract-based Delta Generation}
    \label{alg:subtract-based-delta}
    \begin{algorithmic}[1]
        \Function{SubtractDelta}{$base$, $update$, $PK$}
            \State $\Delta_{new} \gets \pi_{PK}(update \setminus base)$
            \State $\Delta_{removed} \gets \pi_{PK}(base \setminus update) \bowtie \pi_{PK}(update)$
            \State $\Delta_{keys} \gets \Delta_{new} \cup \Delta_{removed}$
            \State \Return $update \bowtie \Delta_{keys}$
        \EndFunction
    \end{algorithmic}
\end{algorithm}

While the subtract-based delta generation worked, it is not efficient. The main bottleneck of the 
OPENIntel system is the retrieval from storage. If both the dataframes do not fit in the allocated RAM,
during operations, they need to be retrieved multiple times for the subtractions after which the update 
dataframe is needed once more for the join with the delta keys. To counter these retrievals, a hash-based 
delta generation algorithm, Algorithm~\ref{alg:hash-based-delta}, has been designed. The hash-based delta
generation uses an auxiliary function to transform a dataframe such that each row is a primary key combined 
with a hash generated from all values in all rows from a block. This way the data stored in memory is smaller 
and only requires one match between dataframes instead of comparing each row with all others like the 
subtract-based generation would. Using this auxiliary function a hash-based dataframe can be generated 
for both the base and the update to then join these dataframes and filter out all keys that have matching 
hashes. Finally, the primary keys that remain will identify all rows from the update to store in the delta.

\begin{algorithm}
    \caption{Hash-based Delta Generation}
    \label{alg:hash-based-delta}
    \begin{algorithmic}[1]

        \Function{HashBlock}{$dataframe$, $PK$}
            \State $dataframe' \gets dataframe$ \textbf{with} $row\_hash \gets H(row)$ \textbf{for each} $row$
            \State \Return $\text{groupBy}(dataframe', PK,\ H(\text{sort}(row\_hashes)))$
        \EndFunction
        \State
        \Function{GenerateHashDelta}{$base$, $update$, $PK$}
            \State $base\_hashes \gets \text{HashBlock}(base, PK)$
            \State $update\_hashes \gets \text{HashBlock}(update, PK)$

            \State $joined \gets update\_hashes \overset{left}{\bowtie} base\_hashes$

            \State $\Delta_{keys} \gets \sigma_{\ base.hash = \emptyset\ \vee\ base.hash \neq update.hash}\ (joined)$

            \State \Return $update \bowtie \Delta_{keys}$
        \EndFunction

    \end{algorithmic}
\end{algorithm}

\subsection{Retrieving Data from the DeltaParq}

Retrieving a dataframe for a given day from the DeltaParq requires exploring the full 
dependency chain, reconstruct all dataframes in the dependency chain, and finally join 
the dataframe with the volatile data of the day that is to be retrieved.

\subsubsection{Dependency Resolution}

Given a delta strategy function $strategy: \mathbb{Z} \rightarrow \mathbb{Z}$ that maps each
day to the day it depends on, the dependency chain for a given day can be resolved for a day by 
walking the chain until a snapshot, identified by $strategy(d) = d$, is reached. If multiple 
days are requested, the dependency chain is walked for each day. If at any point a dependency is 
identified that was already part of the chain of a different day, the chains are merged. This way 
the dependency resolver always returns a chain without duplicates that ensures that each day's 
dependency precedes itself in the chain:
$$\forall d \in chain,\quad strategy(d) = d \;\lor\; strategy(d) \in chain_{[:i_d]}$$
where $chain_{[:i_d]}$ denotes the chain up until $d$.

\subsubsection{Delta Reconstruction}

Using the resolved dependency chain, the full dataframe for each day in the chain can be reconstructed 
which will result in all the requested days being reconstructed as well. For a snapshot day the data is read 
directly from storage, while for a delta day the delta is applied to a previously reconstructed dataframe. 
The application of a delta to a base requires filtering out all primary keys in the delta from the base as 
these have changed between the two measurements using an anti join. Afterward, the result can be unioned with the delta to 
construct the full dataframe:
$$(base \triangleright delta) \cup delta$$
As a final step for reconstructing the full dataframe, the applied delta needs to be joined with the volatile dataframe. 
Even if the volatile columns are not used in the further analysis, the volatile dataframe determines which data was present 
for the specific measurement. Since the delta only stores updates found in the new measurement, blocks that were fully removed 
from in the update cannot be present in the delta dataframe. Since the volatile dataframe measures all the present blocks and 
its volatile columns for each measurement, this can be used to filter out all the completely removed blocks that are still 
present in the applied delta as these would not be present in the volatile dataframe.

\subsubsection{Public Interface}

The DeltaParq offers both a $get()$ and $get\_list()$ interface for retrieving data. With $get\_list()$ the retrieval 
of the dataframes is structured such that any dataframes that would be retrieved multiple times in the different dependency 
chain of each requested day is only retrieved once. Furthermore, the retrieval methods of the DeltaParq offer a column selection 
and filter pushdown. Both the column selection and filter are immediately applied to the dataframes retrieved from storage, like 
pyspark's lazy evaluation would do. This way only the necessary data is retrieved speeding up the retrieval and analysis process.

\subsection{Performance Constraints}

The DeltaParq would benefit significantly from data being physically partitioned by the primary key on write. Both the store 
and retrieval part of the DeltaParq use various joins and substracts which are wide operations. In the case of spark executing 
on multiple excutors, the executors need to share information for these wide operations. An executor may have a row that needs to be joined with a row 
from the other dataframe present on a different executor causing waiting time on other nodes and extra overhead for sharing the data. 
If all data is partitioned such that each partition is guaranteed to have all rows for each primary key it has, no data sharing is 
necessary for these wide operations as each executor now knows they aren't possibly missing data. If this partitioning is done efficiently
on write and pyspark were able to read this, the now broad execution path of combining all the data can be transformed into a direct linear 
path for each executor.

However, while pyspark does support saving partitions in buckets, it is only able to do this by using a Hive metastore. The 
metastore would store all the information on which buckets were used to ensure correctness, the same way DeltaParq uses a stored 
monthly configuration for its self-describing storage. Without a metastore pyspark does not have a guarantee that the data is 
partitioned or bucketed and will always assume no partitioning. Furthermore, pyspark does not offer a way of telling the program 
the data is partitioned in a certain way without the metastore. In the current deployment, no metastore is available, therefore 
the DeltaParq has not been integrated with any form of partitions.


% Keynotes:


% build on top of spark/current used products
% schema decomposition
    % Use config to get the columns and just split that shit open
% delta generation
    % subtraction method
    % "faster" hash method
% delta joining
    % generate dependencies
    % sort dependencies such that base is always build before being used
    % return
% public methods
    % put / put_list
    % get / get_range (explain with/without explicit removal save and select/filter pushdown)
    % rebuild
    
% optimization problems with spark
    % no metastore == no partitioning/order
    % big sad
